\section{Motivations and Backgrounds}

\haochen{I think this section is a bit redundant given that in section 2 we mentioned both the problem with existing works and how our new system tackles the problem. Compared with some recent papers, in KernelBlaster paper they kept only the existing problems in recent work section and left how they resolved the problem in methodology. In TraceLab paper they have one combined "Background and Motivation" section where they mentioned both the current gap in related works and how their method resolved it, which assembles our current outline.}

\haochen{I think it would make more sense to keep this paragraph if we move part of section 2 where we explain how our work is different from previous works to this section. Would love to hear others' ideas. }

\minjia{The common order is intro => background + related work => Challenges/Motivations => Design => Eval. In our case, since there are very few prior work on agentic AD, we don't necessarily need a dedicated challenges/motivation sections if we provide sufficient motivation in the intro. For the challenge part, I expect the benchmark coverage to be the major challenge, e.g.., existing work simply did not fully cover the proper evaluation for AD. Based on that, the order I have in mind would be: intro => background (Triton, AD) + Related work => Benchmark => Design => Eval. One question we may get: what if we apply existing agentic optimization for Triton to Triton-AD, would that already solves most problem, and why do we need another agentic system to solve the Triton-AD? To address that, we do need to compare against some existing solutions that were built to optimize Triton FWD.}
\wangjia{make sense to me.}

\ajay{Agree with Haochen that a standalone section is redundant, and Minjia's proposed order seems more intuitve to me. Per Ahan's suggestion, we can fold motivation into Related Work: I have ablation results on KernelAgent and GEPA showing single-kernel and non-dispatching search frameworks structurally can't co-design the fwd-bwd pair (saved-state contract, peak memory), which concretely motivates why Triton-AD needs a different approach. Can draft a paragraph/ablation once we have more results.}